%

\section{Okręgi ortogonalne, pęki okręgów.}
\index{okrąg!ortogonalny}
% \index{pęk okręgów}
Wie czym są pęki okręgów, zna ich podstawowe własności i potrafi stosować w konfiguracjach spokrewnionych z twierdzeniem Ponceleta.   
\index{twierdzenie!Ponceleta}

Choć pojęcie potęgi punktu względem okręgu można dostrzec już w Propozycjach 35 i 36 Księgi III Elementów Euklidesa, to w pełni zostało ono sformułowane i rozwinięte dopiero przez Louisa Gaultiera w artykule opublikowanym w 1813 roku w Journal de l'École Polytechnique. Właśnie tam po raz pierwszy pojawiają się terminy oś i środek potęgowy, zaś samo określenie „potęga” wprowadził nieco później Jacob Steiner. Badania nad okręgami ortogonalnymi oraz wiązkami współśrodkowymi prowadzili na początku XIX wieku Gaultier, Poncelet, Steiner, J. B. Durrende i inni. Ta stosunkowo nowa geometria elementarna okręgu znalazła cenne zastosowania w różnych działach matematyki i fizyki. Rozpoczniemy od definicji krzywych ortogonalnych. (potęga: zamiana radical ax na power ax itd.)

% TODO: straight line albo circle to stircle, Eves ps. 88/89
Eves \cite[s. 88-90]{eves1_1972} rozpatruje pary prostopadłych okręgów, co prowadzo do definicji potęgi punktu względem okręgu.

O pękach okręgów pisze Delta $\Delta_{24}^{12}$.

% T2.19 tutaj

% Poncelet przeniesiony do czworokątów dwuśrodkowych

%